% Section IV: Chern-Simons theory from anomaly descent and transgression
% Purpose: Motivate CS action physically before diving into structure
% Show how CS emerges from anomaly considerations
%
% Content flow:
% 1. Classical action: S_CS[A] = (k/4π)∫_M tr(A∧dA + 2/3 A∧A∧A)
% 2. Equations of motion: F_A = 0 (flatness!)
% 3. Why this is different from Yang-Mills: no metric needed
% 4. Anomaly descent and transgression identity as physical motivation
%
% Status: REORGANIZE from excellent existing content
% Sources:
% - tex_docs/differential_forms_csimons_foundations_v2.tex (section "From Anomalies to Chern-Simons Theory")
%   Lines 1-450 approximately - anomaly motivation is GOOD
% - tex_docs/anomalies_boundaries_topological_response.tex (anomaly inflow)
%   Lines 1-200 - complement with inflow perspective
% - writings/brian_chern_simons_theory.tex
%   Rigorous transgression identity and Maurer-Cartan form (for mathematical depth)
%
% Action needed: MERGE and REORGANIZE
% - Use v2 as pedagogical base (it has the right physics-first narrative)
% - Supplement with Brian's rigorous proofs where helpful
% - Add differential forms background inline as needed (from foundations.tex)

\section{Chern--Simons Theory from Anomaly Descent}

% TODO: Extract and reorganize from sources above
% Keep physics motivation front and center
% Mathematical rigor secondary but present

% FIGURE FLAG [F-anomaly-to-cs]: Diagram showing:
% 4d anomaly ∫tr(F∧F) → transgression → 3d CS term → boundary anomaly

\subsection{Differential Forms and Gauge Connections}

% TODO: Minimal differential forms background needed for CS action
% Source: differential_forms_csimons_foundations_v2.tex, lines ~1-150
% Keep brief - this is setup for the action

\subsection{Anomaly Descent and the Chern-Simons Form}

% TODO: Main derivation
% Source: differential_forms_csimons_foundations_v2.tex, lines ~150-400
% Show how CS emerges from transgressing ∫tr(F∧F)

\subsection{The Chern--Simons Action}

% TODO: Write down the action explicitly
% S_CS[A] = (k/4π)∫_M tr(A∧dA + 2/3 A∧A∧A)
% Variation and equations of motion: F_A = 0
% Source: multiple files, synthesize

\subsection{Why Chern--Simons Theory is Different}

% TODO: No metric in action (contrast with Yang-Mills)
% No local propagating degrees of freedom
% Flatness as the equation of motion
% This is where "topological" becomes central, not peripheral
